% One editorial system for the five existing Studium collections.
% Load after \documentclass, before book-specific configuration.
\usepackage[utf8]{inputenc}
\usepackage[T1]{fontenc}
\usepackage{lmodern}
\DeclareUnicodeCharacter{2212}{\ensuremath{-}}
\DeclareUnicodeCharacter{2192}{\ensuremath{\to}}
\DeclareUnicodeCharacter{03B1}{\ensuremath{\alpha}}
\usepackage[spanish,es-nodecimaldot]{babel}
\spanishplainpercent
\AtBeginDocument{\spanishdeactivate{"~<>}}
\usepackage[a4paper,margin=2.5cm,headheight=14pt]{geometry}
\usepackage{amsmath,amssymb,mathtools,booktabs,array,longtable,graphicx,xcolor,tikz,enumitem,microtype,titlesec,fancyhdr,listings,etoolbox}
\usetikzlibrary{arrows.meta,calc,positioning,patterns}
\usepackage[breakable,skins]{tcolorbox}
\usepackage[labelfont=bf]{caption}
\usepackage[hyphens]{url}
\usepackage[hidelinks,unicode,hypertexnames=false]{hyperref}
\hypersetup{pdfauthor={Daniel Miguel Tejedor}}
\addto\captionsspanish{\renewcommand{\contentsname}{Índice general}\renewcommand{\tablename}{Cuadro}}
\AtBeginDocument{\normalfont\DeclareFontShape{T1}{lmr}{bx}{sc}{<->ssub*lmr/m/sc}{}}
\definecolor{studiumInk}{HTML}{142C40}
\definecolor{studiumRule}{HTML}{142C40}
\definecolor{studiumBlue}{RGB}{24,74,130}
\definecolor{studiumCyan}{RGB}{0,184,205}
\definecolor{studiumLight}{RGB}{232,240,250}
\definecolor{studiumTipBack}{RGB}{232,244,236}
\definecolor{studiumTipFrame}{RGB}{27,107,58}
\definecolor{studiumDefBack}{RGB}{232,240,250}
\definecolor{studiumDefFrame}{RGB}{24,74,130}
\definecolor{studiumWorkedBack}{RGB}{255,244,230}
\definecolor{studiumWorkedFrame}{RGB}{138,78,12}
\definecolor{studiumCheckBack}{RGB}{243,236,248}
\definecolor{studiumCheckFrame}{RGB}{96,48,130}
\definecolor{studiumCodeBack}{RGB}{245,245,242}
\definecolor{studiumCodeRule}{RGB}{48,48,52}
\titleformat{\chapter}[display]{\normalfont\filright}{\large\mdseries\scshape\chaptertitlename\ \thechapter}{1ex}{\normalfont\huge\bfseries}
\titlespacing*{\chapter}{0pt}{2.5ex plus 1ex minus .2ex}{2.3ex}
\titleformat{name=\chapter,numberless}[display]{\normalfont\filright}{}{0pt}{\normalfont\huge\bfseries}
\titleformat{\section}{\normalfont\Large\bfseries}{\thesection}{.8em}{}
\titleformat{\subsection}{\normalfont\large\bfseries}{\thesubsection}{.7em}{}
\titleformat{\subsubsection}{\normalfont\normalsize\bfseries}{\thesubsubsection}{.7em}{}
\makeatletter
\renewcommand*\l@section{\@dottedtocline{1}{1.5em}{3.5em}}
\@openrightfalse
\let\cleardoublepage\clearpage
\makeatother
\setcounter{secnumdepth}{2}\setcounter{tocdepth}{1}
\setlength{\parindent}{1.5em}\setlength{\parskip}{3pt plus .5pt}
\setlength{\emergencystretch}{2em}\raggedbottom
\setlist{itemsep=2pt,topsep=4pt,leftmargin=*}
\setlist[description]{style=nextline,leftmargin=0pt,labelwidth=0pt,labelindent=0pt}
\allowdisplaybreaks[2]
\newcounter{studiumfigure}
\makeatletter\@addtoreset{studiumfigure}{chapter}\makeatother
\renewcommand{\thestudiumfigure}{\thechapter.\arabic{studiumfigure}}
\newcommand{\studiumcaption}[2]{\refstepcounter{studiumfigure}\label{#1}\addcontentsline{lof}{figure}{\protect\numberline{\thestudiumfigure}#2}{\small\textbf{Figura~\thestudiumfigure.} #2\par}}
\newcommand{\Raire}{287.05287}
\newcommand{\gzero}{9.80665}
\tcbset{studiumcommon/.style={breakable,boxrule=.65pt,arc=.7pt,left=1.7mm,right=1.7mm,top=1.3mm,bottom=1.3mm,before skip=10pt,after skip=10pt,fonttitle=\bfseries,coltitle=white,coltext=black},
 studiumtip/.style={studiumcommon,colback=studiumTipBack,colframe=studiumTipFrame,colbacktitle=studiumTipFrame},
 studiumdef/.style={studiumcommon,colback=studiumDefBack,colframe=studiumDefFrame,colbacktitle=studiumDefFrame},
 studiumworked/.style={studiumcommon,colback=studiumWorkedBack,colframe=studiumWorkedFrame,colbacktitle=studiumWorkedFrame},
 studiumcheck/.style={studiumcommon,colback=studiumCheckBack,colframe=studiumCheckFrame,colbacktitle=studiumCheckFrame}}
\newtcolorbox{studiumworked}[1]{studiumworked,title={#1}}
\newcommand{\studiumchapterguide}[1]{\par\noindent\textcolor{studiumInk}{\textbf{Orientación del capítulo.}}\enspace\textit{#1}\par\medskip}
\newcommand{\objetivos}[1]{\studiumchapterguide{#1}}
\newcommand{\caja}[2]{\begin{tcolorbox}[studiumdef,title={#1}]#2\end{tcolorbox}}
\newcommand{\resumen}[1]{\begin{tcolorbox}[studiumcheck,title={Síntesis del capítulo}]#1\end{tcolorbox}}
\newcommand{\studiumchaptercheck}{\begin{tcolorbox}[studiumcheck,title={Comprobación final}]
\begin{enumerate}
\item Explica el modelo del capítulo y las hipótesis que permiten utilizarlo.
\item Reproduce un desarrollo o un ejemplo sin consultar su solución y comprueba unidades, signos y condiciones de validez.
\item Resuelve las preguntas o ejercicios disponibles y distingue el resultado calculado de su interpretación física.
\end{enumerate}\end{tcolorbox}}
\newcommand{\figura}[2]{\begin{figure}[tbp]\centering\includegraphics[width=.83\linewidth]{#1.pdf}\caption{#2}\label{fig:#1}\end{figure}\noindent La figura~\ref{fig:#1} ofrece una representación del modelo discutido.}
\newsavebox{\studiumMathBox}
\newcommand{\fitmath}[1]{\sbox{\studiumMathBox}{$\displaystyle #1$}\ifdim\wd\studiumMathBox>\linewidth\resizebox{\linewidth}{!}{\usebox{\studiumMathBox}}\else\usebox{\studiumMathBox}\fi}
\pagestyle{fancy}\fancyhf{}
\fancyhead[LE]{\footnotesize\thepage\hspace{1.5em}\mdseries\scshape\nouppercase{\leftmark}}
\fancyhead[RO]{\footnotesize\mdseries\scshape\nouppercase{\leftmark}\hspace{1.5em}\thepage}
\renewcommand{\headrulewidth}{.4pt}\renewcommand{\footrulewidth}{0pt}
\fancypagestyle{plain}{\fancyhf{}\fancyfoot[C]{\small\thepage}\renewcommand{\headrulewidth}{0pt}}
\lstdefinestyle{studium}{basicstyle=\ttfamily\small,columns=fullflexible,keepspaces=true,showstringspaces=false,breaklines=true,breakatwhitespace=false,frame=leftline,framerule=1.15pt,rulecolor=\color{studiumCodeRule},backgroundcolor=\color{studiumCodeBack},xleftmargin=1.15em,framexleftmargin=.75em,aboveskip=.9em,belowskip=.7em,tabsize=4,keywordstyle=\ttfamily\bfseries,commentstyle=\ttfamily\itshape,stringstyle=\ttfamily}
\lstset{style=studium}
\input{../../../shared/latex/studium-cover.tex}
\providecommand{\studiumvolumetype}{manual}
\newcommand{\studiumreadingguide}{%
\chapter*{Cómo utilizar este volumen}\addcontentsline{toc}{chapter}{Cómo utilizar este volumen}
\ifdefstring{\studiumvolumetype}{manual}{\noindent\textbf{Ruta de estudio.} Empieza por la orientación del capítulo y sus conceptos. Escribe las hipótesis antes de utilizar las ecuaciones y sigue los desarrollos con sus unidades. Después intenta los problemas y las preguntas de repaso sin consultar las soluciones. Cierra cada capítulo con la comprobación final.}{%
\ifdefstring{\studiumvolumetype}{problemas}{\noindent\textbf{Ruta de trabajo.} Lee el enunciado y copia sus datos con unidades. Define las incógnitas, el modelo y las hipótesis. Intenta resolver el problema antes de consultar el desarrollo. Compara después el resultado, la interpretación física y las comprobaciones.}{\noindent\textbf{Ruta de entrenamiento.} Lee las instrucciones particulares de cada simulacro y reserva el tiempo propuesto. Trabaja primero con los enunciados. Después utiliza las soluciones indicadas en el volumen para revisar el planteamiento, los cálculos y las comprobaciones.}}

\noindent\textbf{Recuadros.} El azul identifica definiciones y resultados; el verde, consejos; el naranja, problemas resueltos; y el violeta, síntesis y autoevaluación. Los apartados técnicos conservan los nombres propios de cada asignatura.

\noindent\textbf{Problemas y simulacros.} Separa enunciado, datos, hipótesis, modelo, desarrollo, resultado y comprobación. En los simulacros trabaja primero con los enunciados y consulta después las soluciones. Sus tiempos y puntuaciones son propuestas de entrenamiento cuando así se indican, sin sustituir las instrucciones oficiales del curso.

\noindent\textbf{Consulta.} El índice general permite localizar los temas. La notación, los formularios, las soluciones y la bibliografía mantienen sus apartados al final del manual cuando forman parte de él.
\section*{Créditos y alcance editorial}\addcontentsline{toc}{section}{Créditos y alcance editorial}
\noindent Daniel Miguel Tejedor. Colección Studium. Edición de estudio independiente para el Grado en Ingeniería Aeroespacial, Universidad de León.

\noindent La compilación y la unificación editorial no añaden una aprobación institucional ni una nueva verificación científica. Cada libro conserva las hipótesis, límites y comprobaciones de sus modelos; los informes registran su alcance y sus pendientes.

\noindent\textbf{Imagen de portada.} \studiumphotocredit
}
